\documentclass[12pt, polish, aspectratio = 169]{slides}

\title{Rozprawa Doktorska}
\author{Łukasz Dróżdż}

\subject{Analiza metrologiczna algorytmów dyskretnej transformacji falkowej}
\subtitle{Analiza metrologiczna algorytmów dyskretnej transformacji falkowej}
\institute{Politechnika Śląska, Wydział Elektryczny \\ Katedra Metrologii, Elektroniki i Automatyki}
\keywords{dyskretna transformacja falkowa, cyfrowe przetwarzanie sygnałów, szacowanie niepewności wyniku pomiaru, analiza właściwości metrologicznych toru pomiarowego}

\titlegraphic{\includegraphics[height = 1cm]{obrazki/polsl_logo}}
\date{30 października 2024}

\def\fullauthor{mgr inż. Łukasz Dróżdż}
\def\fullpromoter{dr hab. inż. Jerzy Roj, prof. PŚ}
\def\fullcontact{+48 32 237 20 05 \\ lukasz.drozdz@polsl.pl}

\begin{document}

\section*{Wstęp do prezentacji oraz plan wystąpienia}

\begin{frame}[plain]
\titlepage
\end{frame}

\begin{frame}{Plan prezentacji}
\tableofcontents
\end{frame}

\newsection{Motywacja pracy}{Przedstawienie motywacji pracy i aktualnego stanu wiedzy}

\begin{frame}{Schemat blokowy toru pomiarowego}
\begin{figure}
\includegraphics[scale = 0.85]{obrazki/schemat_demo}
\caption{Schemat blokowy toru pomiarowego zawierającego w swojej strukturze algorytm transformacji falkowej}
\end{figure}
\begin{itemize}
\item $s(t)$ -- zmienna w czasie wielkość fizyczna.
\item $y(t)$ -- napięciowa reprezentacja wielkości $s(t)$.
\item $x(i)$ -- kolejne próbki wielkości $y(t)$.
\item $\mathbfit{X}(i)$ -- wektor wielkości wyjściowych algorytmu.
\end{itemize}
\end{frame}

\begin{frame}{Algorytm transformacji falkowej}
\begin{equation}
w_{a,b} = \frac{1}{\sqrt{a}} \int _{-\infty} ^{\infty} x \emb{t} \psi \emb{\frac{t-b}{a}} dt \label{eq:cwt}
\end{equation}
\begin{itemize}
\item Umożliwia analizę sygnału $x(t)$ w dziedzinie skala-czas (informacja o czasie).
\item Wykorzystuje funkcję \enquote{falki-matki}, oznaczaną jako $\psi(t)$.
\item Parametr skali $a$ określa wyodrębniany zakres częstotliwości.
\item Parametr $b$ określa przesunięcie falki w czasie (okno pomiarowe).
\item Współczynniki $w_{a,b}$ stanowią wielkości wyjściowe algorytmu.
\end{itemize}
\begin{itemize}
\item[CWT] ciągła transformacja falkowa (w praktyce niemożliwa do realizacji),
\item[DWT] dyskretna transformacja falkowa (dla $a = 2^m$, $b = n2^m$, gdzie $m, n \in \mathbb{N}$),
\item[UFWT] dyskretna transformacja falkowa bez decymacji (bez odrzucania próbek).
\end{itemize}
\end{frame}

\begin{frame}{Przykłady zastosowań algorytmu}
\begin{itemize}
\item Przetwarzanie, kompresja oraz analiza obrazów i dźwięków.
\item Pomiary termowizyjne.
\item Redukcja szumu w sygnale pomiarowym.
\item Analiza sygnałów EKG oraz EEG.
\item Analiza przebiegu drgań sejsmicznych.
\item Bezinwazyjna analiza stanu elementów mechanicznych.
\item Algorytmy stratnej kompresji danych.
\item Analiza zwarć wysoko prądowych.
\item Analiza procesów hydrologicznych.
\item Wiele innych zastosowań oraz ciągły ich rozwój.
\end{itemize}
\end{frame}

\begin{frame}{Problem analizy właściwości metrologicznych}
\begin{itemize}
\item Możliwość parametryzacji stosowanego algorytmu:
	\begin{itemize}
	\item zmiana rodzaju stosowanej falki-matki,
	\item zmiana liczby iteracji procesu dekompozycji.
	\end{itemize}
\item Równanie falki matki $\psi(t)$ nie zawsze jest dostępne:
	\begin{itemize}
	\item wyprowadzenie odpowiednich zależności jest czasochłonne,
	\item wymagana jest ekspercka wiedza na temat stosowanej rodziny falek.
	\end{itemize}
\item Zmiana parametrów stosowanego algorytmu wymusza ponowną analizę:
	\begin{itemize}
	\item metody analityczne są skomplikowane i nie są uniwersalne,
	\item metoda Monte Carlo jest czasochłonna.
	\end{itemize}
\item Brak wskazówek w Przewodniku JCGM:
	\begin{itemize}
	\item analiza właściwości metrologicznych jest pomijana.
	\end{itemize}
\end{itemize}
\end{frame}

\newsection{Teza pracy}{Przedstawienie tezy pracy oraz jej najważniejszych założeń}

\begin{frame}{Teza pracy}
\begin{justify}
Stosując przedstawiony w pracy model błędów oraz zaproponowaną metodę szacowania wartości wypadkowej niepewności rozszerzonej istnieje możliwość oszacowania wartości niepewności rozszerzonych dla wielkości wyjściowych toru pomiarowego wykorzystującego algorytm dyskretnej transformacji falkowej.

Oszacowanie wartości niepewności rozszerzonych dla omawianych wielkości jest możliwe w trakcie działania systemu pomiarowego, również w przypadku zmiany parametrów pracy tego systemu oraz zmiany parametrów modelu błędów.
\end{justify}
\end{frame}

\newsection{Model błędów}{Omówienie najważniejszych właściwości zaproponowanego modelu błędów}

\begin{frame}{Model błędów toru pomiarowego}
\begin{figure}
\includegraphics[scale = 0.85]{obrazki/schemat_trans}
\caption{Model błędów pojedynczej wielkości wyjściowej toru pomiarowego}
\end{figure}
\begin{itemize}
\item $s(t)$ -- zmienna w czasie wielkość fizyczna.
\item $y(t)$ -- napięciowa reprezentacja wielkości $s(t)$.
\item $x(i)$ -- kolejne próbki wielkości $y(t)$.
\item $X(i)$ -- wielkość wyjściowa toru pomiarowego.
\item $G_{y}(j\omega)$, $H_{X}(z)$ -- transmitancja części analogowej/cyfrowej.
\item $f_{y}(x)$, $f_{X}(x)$ -- funkcja przetwarzania części analogowej/cyfrowej.
\end{itemize}
\end{frame}

\begin{frame}{Algorytm transformacji falkowej}
\begin{equation}
\begin{bmatrix}
X _{0}   \\
X _{1}   \\
\vdots      \\
X _{M-1}
\end{bmatrix}
=
\begin{bmatrix}
a_{0, 0}   &   a_{0, 1} &   \cdots   &   a_{0, N-1}      \\
a_{1, 0}   &   \ddots   &            &   a_{1, N-1}      \\
\vdots     &            &   \ddots   &   \vdots          \\
a_{M-1, 0} &   \cdots   &   \cdots   &   a_{M-1, N-1}
\end{bmatrix}
\begin{bmatrix}
x _{0}   \\
x _{1}   \\
\vdots      \\
x _{N-1}
\end{bmatrix}
\label{eq:alg_out_mat}
\end{equation}
\begin{itemize}
\item $X_{i}$ -- wielkości wyjściowe algorytmu ($M$ wielkości).
\item $x_{j}$ -- wielkości wejściowe algorytmu ($N$ wielkości).
\item $a_{i,j}$ -- współczynniki macierzy transformacji algorytmu:
	\begin{itemize}
	\item wyznaczane analitycznie, na podstawie właściwości falki-matki,
	\item identyfikowane stosując gotową implementację algorytmu.
	\end{itemize}
\end{itemize}
\end{frame}

\begin{frame}{Właściwości algorytmu transformacji falkowej}
\begin{columns}
\begin{column}{0.5\textwidth}
	\begin{figure}
	\includegraphics[scale = 0.85]{obrazki/dwt_trans}
	\caption{Model właściwości dynamicznych algorytmu transformacji falkowej}
	\end{figure}
\end{column}
\begin{column}{0.5\textwidth}
	\begin{itemize}
	\item $T_{j}(i)$, $S_{j}(i)$ -- detale oraz aproksymacje dla $j$-tej skali.
	\item $x(i)$ -- wielkości wejściowe algorytmu.
	\item $N_{d}$ -- liczba iteracji procesu dekompozycji sygnału wejściowego.
	\item $H_{*}(z)$ -- transmitancja obiektu związanego z wybraną grupą wielkości wyjściowych.
	\end{itemize}
\end{column}
\end{columns}
\end{frame}

\begin{frame}{Geneza sygnałów błędów własnych}
\begin{enumerate}
\item \label{item:trans_err_group} Wynikające z niedoskonałości transmitancji algorytmu:
	\begin{itemize}
	\item istnieje możliwość deterministycznego opisu przebiegu sygnałów błędów,
	\item projektant toru pomiarowego definiuje postać transmitancji $\dot{H}_{i}(z)$,
	\item sygnały te występują gdy $\tilde{H}_{i}(z) \ne \dot{H}_{i}(z)$.
	\end{itemize}
\item \label{item:numeric_err_group} Wynikające z rzeczywistej implementacji algorytmu:
	\begin{itemize}
	\item genezą są ograniczone możliwości precyzji zapisu liczb zmiennoprzecinkowych,
	\item możliwość opisu parametrów w kategoriach probabilistycznych.
	\end{itemize}
\item Wynikające z niedoskonałości wyznaczenia wartości współczynników macierzy transformacji:
	\begin{itemize}
	\item dla błędów numerycznych można zaliczyć je do grupy~\itemref{item:numeric_err_group},
	\item w pozostałych przypadkach można zaliczyć je do grupy~\itemref{item:trans_err_group}.
	\end{itemize}
\end{enumerate}
\end{frame}

\begin{frame}{Parametry sygnałów błędów własnych}
\begin{columns}
\begin{column}{0.6\textwidth}
	\begin{equation}
	e_{X_{j}}\emb{i} = \tilde{X}_{j}\emb{i} - \dot{X}_{j}\emb{i} \label{eq:dwt_round_error}
	\end{equation}
	\begin{figure}
	\includegraphics[scale = 0.85]{obrazki/schemat_dwt_ew_min}
	\caption{Proces identyfikacji parametrów sygnałów błędów własnych dla $j$-tej wielkości wyjściowej algorytmu}
	\end{figure}
\end{column}
\begin{column}{0.4\textwidth}
	\begin{itemize}
	\item Parametry sygnału $e_{X_{j}}(i)$ wyznaczane są stosując rzeczywistą implementacje algorytmu.
	\item Sygnały błędów wielkości wejściowych $x(i)$ nie występują.
	\item Parametry sygnału wejściowego $x(i)$ muszą być zbliżone do rzeczywistych.
	\end{itemize}
\end{column}
\end{columns}
\end{frame}

\newsection{Niepewność rozszerzona}{Opis metody wyznaczania wartości wypadkowej niepewności rozszerzonej}

\begin{frame}{Miara niedokładności wyniku pomiaru}
\begin{columns}
\begin{column}{0.5\textwidth}
	\begin{figure}
	\includegraphics[scale = 0.85]{obrazki/hist_test}
	\caption{Graficzna interpretacja miary niepewności rozszerzonej jako parametru sygnału błędu $e(t)$}
	\end{figure}
\end{column}
\begin{column}{0.5\textwidth}
	\begin{equation}
	P \emb{\left|\hat{e} \emb{t}\right| \le U} = \gamma \label{eq:unc_definition}
	\end{equation}
	\begin{itemize}
	\item Wyznaczana dla wybranego poziomu ufności $\gamma = 1 - \alpha$.
	\item Dostarcza więcej informacji, niż niepewność standardowa.
	\item Wyznaczenie wartości wypadkowej niepewności rozszerzonej może być skomplikowane.
	\end{itemize}
\end{column}
\end{columns}
\end{frame}

\begin{frame}{Wypadkowa niepewność rozszerzona}
\begin{itemize}
\item Analizowana wielkość $s(t)$ jest zakłócona wieloma sygnałami błędów, przy czym $\tilde{s}(t) = \dot{s}(t) + e_{s,\Sigma}(t)$ oraz $e_{s,\Sigma}(t) = e_{s,0}(t) + \hdots + e_{s,N-1}(t)$.
\item Rozważane przypadki dla zmiennej liczby i parametrów sygnałów błędów:
	\begin{itemize}
	\item warunki centralnego twierdzenia granicznego nie zawsze są spełnione.
	\end{itemize}
\item Właściwości algorytmu transformacji falkowej:
	\begin{itemize}
	\item wzmocnia/tłumi sygnały błędów o określonym widmie,
	\item osobna analiza dla każdej grupy wielkości wyjściowych.
	\end{itemize}
\item Problemy:
	\begin{itemize}
	\item metoda Monte Carlo jest przystępna, jednak bardzo czasochłonna,
	\item metody analityczne są skomplikowane i wymagają wyprowadzeń.
	\end{itemize}
\end{itemize}
\end{frame}

\begin{frame}{Metoda redukcyjnej arytmetyki interwałowej}
\begin{equation}
U_{\Sigma} = \sqrt{
\begin{bmatrix}
U_{0} \\ U_{1} \\ \vdots \\ U_{N-1}
\end{bmatrix}^{T}
\begin{bmatrix}
1         & h_{0,1} & \cdots & h_{0,N-1} \\
h_{1,0}   & 1       &        & h_{1,N-1} \\
\vdots    &         & \ddots & \vdots    \\
h_{N-1,0} & \cdots  & \cdots & 1
\end{bmatrix}
\begin{bmatrix}
U_{0} \\ U_{1} \\ \vdots \\ U_{N-1}
\end{bmatrix}}
\label{eq:unc_matrix}
\end{equation}
\begin{itemize}
\item $U_{\Sigma}$ -- oszacowana wartość wypadkowej niepewności rozszerzonej.
\item $U_{i}$ -- wartość niepewności rozszerzonej $i$-tego sygnału błędu.
\item $h_{i,j} = h_{j,i}$ -- współczynnik koherencji dla pary sygnałów $e_{i}(t)$ oraz $e_{j}(t)$:
	\begin{itemize}
	\item wartości zależą od parametrów analizowanych sygnałów błędów,
	\item wyznaczane stosując metodę analityczną lub symulacyjną,
	\item szacowane stosując odpowiednią metodę.
	\end{itemize}
\end{itemize}
\end{frame}

\begin{frame}{Aplikacja omawianej metody}
\begin{itemize}
\item Wartości współczynników kształtu wyznaczane na podstawie równania:
\begin{equation}
s_{a,b} = s_{b,a} = \frac{U_{a,b}^{2} - U_{a}^{2} - U_{b}^{2}}{2 U_{a} U_{b}} = \frac{U^{2}}{2 U_{a,b}} - 1 \label{eq:unc_shapertwo}
\end{equation}
\item Korekta wynikająca z centralnego twierdzenia granicznego:
\begin{equation}
k_{i,j} = k_{j,i} = \frac{U_{i}^{2} + U_{j}^{2}}{\sum _{k = 0} ^{N-1} U_{k}^{2}} \label{eq:unc_cohercorrb}
\end{equation}
\item Korekta wynikająca z dysproporcji wartości niepewności rozszerzonych:
\begin{equation}
p_{a,b} = p_{b,a} = \sqrt{\frac{\min \emb{U_{a}, U_{b}}}{\max \emb{U_{a}, U_{b}}}} \label{eq:unc_cohercorra}
\end{equation}
\end{itemize}
\end{frame}

\begin{frame}{Aplikacja omawianej metody}
\begin{itemize}
\item Wartości współczynników koherencji szacowane zgodnie z równaniem:
\begin{equation}
h_{i,j} = h_{j,i} = s_{i,j} p_{i,j} k_{i,j} = s_{i,j} \sqrt{\frac{\min \emb{U_{i}, U_{j}}}{\max \emb{U_{i}, U_{j}}}} \left( \frac{U_{i}^{2} + U_{j}^{2}}{\sum _{k = 0} ^{N-1} U_{k}^{2}} \right) \label{eq:unc_coher}
\end{equation}
\item W celu aplikacji metody należy:
	\begin{itemize}
	\item wyznaczyć wartości współczynników kształtu -- operacja jednorazowa,
	\item skorygować wartości współczynników kształtu dla obecnych parametrów.
	\end{itemize}
\item Najistotniejsze cechy metody:
	\begin{itemize}
	\item niska złożoność obliczeniowa -- możliwość stosowania w czasie rzeczywistym,
	\item możliwość zmiany liczby oraz parametrów sygnałów błędów,
	\item względny błąd oszacowania wartości $U_{\Sigma}$ mieszczący się w przedziale \qty{\pm 5}{\percent}.
	\end{itemize}
\end{itemize}
\end{frame}

\newsection{Weryfikacja tezy}{Przedstawienie wybranych wyników badań weryfikujących tezę pracy}

\begin{frame}{Weryfikacja założeń przedstawionych w tezie}
\begin{itemize}
\item Metoda szacowania wartości wypadkowej niepewności rozszerzonej:
	\begin{itemize}
	\item badania symulacyjne przeprowadzone metodą Monte Carlo, 
	\item uwzględnienie zmiennej liczby oraz zmiennych parametrów sygnałów błędów,
	\item weryfikacja wpływu zaproponowanej korekty na dokładność wyników,
	\end{itemize}
\item Zaproponowany w pracy model błędów:
	\begin{itemize}
	\item weryfikacja symulacyjna metodą Monte Carlo,
	\item weryfikacja pomiarowa na przykładzie prototypu toru pomiarowego.
	\end{itemize}
\end{itemize}
\end{frame}

\begin{frame}{Wartość wypadkowej niepewności rozszerzonej}
\begin{figure}
\includegraphics[scale = 0.85]{obrazki/hist_rederr_all_2_5_1_20}
\caption{Względny błąd oszacowania wartości wypadkowej niepewności rozszerzonej \textbf{a)}~stosując proponowaną metodę, \textbf{b)}~stosując oryginalną metodę (bez korekty), \textbf{c)}~stosując Przewodnik JCGM, $N \in \interval{2}{5}$, $U_{i} \in \interval{1}{20}$, $1-\alpha = \qty{95}{\percent}$}
\end{figure}
\end{frame}

\begin{frame}{Eksperyment symulacyjny}
\begin{figure}
\includegraphics[scale = 0.85]{obrazki/schemat_symul}
\caption{Schemat blokowy toru pomiarowego będącego obiektem eksperymentu symulacyjnego}
\end{figure}
\begin{itemize}
\item $s(t)$ -- zmienna w czasie wielkość fizyczna.
\item $y_{*}(t)$ -- napięciowa reprezentacja wielkości $s(t)$.
\item $x_{*}(i)$ -- dyskretne próbki napięcia reprezentującego wielkość $s(t)$.
\item $\mathbfit{X}(k)$ -- wielkości wyjściowe dla $k$-tej realizacji pomiaru.
\end{itemize}
\end{frame}

\begin{frame}{Wyniki eksperymentu symulacyjnego}
\begin{table}
\small\caption{Zestawienie wyników eksperymentu symulacyjnego \textbf{a)}~stosując proponowaną metodę, \textbf{b)}~stosując oryginalną metodę (bez korekty), \textbf{c)}~stosując Przewodnik JCGM, \textbf{s)}~stosując symulację metodą Monte Carlo (średnia dotyczy wszystkich wielkości)}
\begin{tabular}[c]{| c *{4}{|S[table-format = 3.3]} *{3}{|S[table-format = +1.2]} |} \hline
\multirow{2}{*}{\textbf{Wielkość}} & \multicolumn{4}{c|}{\textbf{Niepewność, mV}} & \multicolumn{3}{c|}{\textbf{Błąd, \%}} \\ \cline{2-8}
          & $U_{c}$ & $U_{a}$ & $U_{b}$ & $U_{s}$ & $\delta_{c}$ & $\delta_{a}$ & $\delta_{b}$ \\ \hline
$S_{2,0}$ & 75.01 & 74.10 & 77.32 & 72.87 & +2.94 & +1.69 & +6.11 \\ \hline
$T_{2,0}$ & 57.29 & 55.91 & 57.58 & 53.89 & +6.31 & +3.75 & +6.85 \\ \hline
$T_{1,0}$ & 47.31 & 43.37 & 46.09 & 43.76 & +8.11 & -0.89 & +5.32 \\ \hline
$T_{1,3}$ & 44.79 & 43.39 & 45.43 & 43.17 & +3.75 & +0.51 & +5.24 \\ \hline
\hline\multicolumn{5}{|c|}{Średnia wartości bezwzględnych wielkości $\delta_{*}$} & 5.77 & 1.48 & 6.13 \\ \hline
\end{tabular}
\end{table}
\end{frame}

\begin{frame}{Eksperyment pomiarowy}
\begin{figure}
\includegraphics[scale = 0.85]{obrazki/schemat_real}
\caption{Schemat blokowy prototypu toru pomiarowego}
\end{figure}
\begin{itemize}
\item $s(t)$ -- zmienne w czasie napięcie wejściowe (RIGOL DG1011).
\item $x(i)$ -- wielkości wejściowe algorytmu transformacji falkowej.
\item $\mathbfit{X}(j)$ -- wielkości wyjściowe dla $j$-tej realizacji pomiaru.
\end{itemize}
\end{frame}

\begin{frame}{Aplikacja zaproponowanej metody analizy}
\begin{enumerate}
\item Określenie, w jaki sposób fragmenty propagują sygnały błędów.
\item Identyfikacja źródeł sygnałów błędów własnych:
	\begin{itemize}
	\item zawartych w sygnale $s(t)$ (parametrów generatora),
	\item wzmacniacza (właściwości statyczne i dynamiczne),
	\item przetwornika A/C (kwantowanie, dynamika układu P/P),
	\item algorytmu transformacji falkowej (operacje zaokrągleń).
	\end{itemize}
\item Identyfikacja parametrów modelu błędów:
	\begin{itemize}
	\item określenie wariancji, wartości oczekiwanej oraz niepewności rozszerzonej,
	\item określenie charakteru realizacji oraz widma sygnałów błędów,
	\item określenie parametrów rozkładu realizacji sygnałów błędów,
	\item określenie korelacji analizowanych sygnałów błędów.
	\end{itemize}
\item Aplikacja zaproponowanej metody szacowania wartości wypadkowej niepewności rozszerzonej.
\end{enumerate}
\end{frame}

\begin{frame}{Weryfikacja wartości niepewności wielkości $x(i)$}
\begin{table}[p]
\small\caption{Zmierzone oraz oszacowane wartości wypadkowej niepewności rozszerzonej wielkości $x(i)$ (znane parametry sygnału $s(t)$, średnia dla całego eksperymentu)}
\begin{tabular}[c]{| c *{2}{|S[table-format = 4.3]} *{2}{|S[table-format = 4.4]} | S[table-format = +2.2] |} \hline
\multirow{2}{*}{\textbf{$f_{s,o}$, Hz}} & \multicolumn{2}{c|}{\textbf{Wariancja, \micro V\textsuperscript{2}}} & \multicolumn{2}{c|}{\textbf{Niepewność, mV}} & \textbf{Błąd, \%} \\ \cline{2-6}
& $\sigma_{c}^{2}$ & $\sigma_{m}^{2}$ & $U_{c}$ & $U_{m}$ & $\delta_{c}$ \\ \hline
600     &       0.59    &       0.44    &       1.24    &       1.16    &       +6.74    \\ \hline
4000    &       24.71   &       20.51   &       7.06    &       6.57    &       +7.48    \\ \hline
8000    &       104.47  &       105.94  &       14.44   &       14.61   &       -1.14    \\ \hline
14000   &       347.53  &       337.12  &       26.31   &       25.84   &       +1.81    \\ \hline
21000   &       857.76  &       775.38  &       41.31   &       39.16   &       +5.49    \\ \hline
\hline\multicolumn{5}{|c|}{Średnia wartości bezwzględnych wielkości $\delta_{*}$} & 5.21 \\ \hline
\end{tabular}
\end{table}
\end{frame}

\begin{frame}{Weryfikacja wartości niepewności wielkości $X_{24}(i)$}
\begin{table}[p]
\small\caption{Zmierzone oraz oszacowane wartości wypadkowej niepewności rozszerzonej wielkości $X_{24}(i)$ (znane parametry sygnału $s(t)$, średnia dla całego eksperymentu)}
\begin{tabular}[c]{| c *{2}{|S[table-format = 4.3]} *{2}{|S[table-format = 4.4]} | S[table-format = +2.2] |} \hline
\multirow{2}{*}{\textbf{$f_{s,o}$, Hz}} & \multicolumn{2}{c|}{\textbf{Wariancja, \micro V\textsuperscript{2}}} & \multicolumn{2}{c|}{\textbf{Niepewność, mV}} & \textbf{Błąd, \%} \\ \cline{2-6}
& $\sigma_{c}^{2}$ & $\sigma_{m}^{2}$ & $U_{c}$ & $U_{m}$ & $\delta_{c}$ \\ \hline
600     &       0.51    &       0.48    &       1.39    &       1.37    &       +1.90     \\ \hline
4000    &       224.98  &       186.66  &       21.28   &       19.68   &       +8.12     \\ \hline
8000    &       920.05  &       932.81  &       42.85   &       43.28   &       -0.97     \\ \hline
14000   &       1.07    &       1.88    &       1.96    &       2.73    &       -28.29    \\ \hline
21000   &       2.19    &       32.09   &       2.57    &       8.75    &       -70.59    \\ \hline
\hline\multicolumn{5}{|c|}{Średnia wartości bezwzględnych wielkości $\delta_{*}$} & 12.69 \\ \hline
\end{tabular}
\end{table}
\end{frame}

\begin{frame}{Przyczyny rozbieżności uzyskanych wyników}
\begin{figure}
\includegraphics[scale = 0.85]{obrazki/schemat_selektywnosc}
\caption{Stosunek wariancji sygnałów błędów dla wybranego wariantu eksperymentu}
\end{figure}
\end{frame}

\newsection{Podsumowanie}{Przedstawienie najważniejszych osiągnięć i wniosków płynących z pracy}

\begin{frame}{Najważniejsze osiągnięcia}
\begin{itemize}
\item Zaproponowanie jednolitego modelu błędów, opisującego właściwości toru pomiarowego, zarówno części cyfrowej oraz analogowej.
\item Przedstawienie algorytmu transformacji falkowej w postaci macierzowej wraz z metodami identyfikacji wartości współczynników macierzy transformacji.
\item Wskazanie, w jaki sposób algorytm transformacji falkowej przetwarza sygnały błędów z wejścia na wyjście.
\item Opis sygnałów błędów własnych algorytmu transformacji falkowej oraz sposobu identyfikacji ich parametrów.
\item Zaproponowanie przystępnej w aplikacji metody szacowania wartości wypadkowej niepewności rozszerzonej, możliwej do stosowania w czasie rzeczywistym dla zmiennych parametrów modelu błędów.
\item Symulacyjna oraz pomiarowa weryfikacja przedstawionych w pracy zależności oraz wskazanie przykładu aplikacji proponowanej metody analizy.
\end{itemize}
\end{frame}

\begin{frame}{Wnioski}
\begin{itemize}
\item Zaproponowany model błędów jest zasadny w przypadku analizy torów pomiarowych wykorzystujących algorytmy transformacji falkowej, natomiast dokładność oszacowania wartości wypadkowej niepewności rozszerzonej zależy od dokładności określenia parametrów tego modelu.
\item Zaproponowana metoda szacowania wartości wypadkowej niepewności rozszerzonej zapewnia wyniki zbliżone do uzyskiwanych metodą Monte Carlo, a niska złożoność obliczeniowa umożliwia jej aplikację w czasie rzeczywistym, również w przypadku zmiany parametrów modelu błędów.
\item Stosowanie wskazanej metody możliwe jest zawsze, jednak uzasadnione jest wtedy, gdy warunki centralnego twierdzenia granicznego nie są spełnione.
\item Mimo spełnienia warunków centralnego twierdzenia granicznego dla wielkości wejściowych, warunki te mogą nie być spełnione dla wielkości wyjściowych algorytmu transformacji falkowej.
\end{itemize}
\end{frame}

\begin{frame}{Dorobek naukowy}
\begin{itemize}
\item 13 współautorskich artykułów naukowych w czasopismach:
	\begin{itemize}
	\item[6:] Przegląd Elektrotechniczny (IF 0,5),
	\item[3:] Applied Sciences-Basel (IF 2,5),
	\item[1:] Electronics (Switzerland) (IF 2,69),
	\item[1:] International Journal of Electronics and Telecommunications (IF 0,5),
	\item[2:] Zesz. Nauk. Wydziału Elektrotechniki i Automatyki Politechniki Gdańskiej.
	\end{itemize}
\item 9 referatów wygłoszonych na konferencjach krajowych:
	\begin{itemize}
	\item[3:] Podstawowe Problemy Metrologii (2024, 2023, 2022),
	\item[3:] Systemy pomiarowe w badaniach naukowych i przemyśle (2024, 2022, 2018),
	\item[3:] Międzyuczelniana Konferencja Metrologów (2024, 2020, 2019).
	\end{itemize}
\item 6 rozdziałów w monografiach współautorskich.
\end{itemize}
\end{frame}

\section*{Zakończenie prezentacji i podziękowania}

\begin{frame}[plain]
\lastpage
\end{frame}

\end{document}
